Understanding how large language models (LLMs) internally represent different types of beliefs is critical for AI safety and hallucination detection.
While prior work has studied truth and knowledge encoding in LLM layers, no study has systematically investigated whether LLMs distinguish between \emph{epistemic} beliefs (evidence-based factual claims) and \emph{non-epistemic} beliefs (opinions, values, and preferences).
We conduct a layer-wise probing analysis across four GPT-2 models (117M--1.5B parameters), extracting residual stream activations and applying mass-mean probing and logistic regression to classify belief types and their truth/agreement values.
We find three key results: (1) all models achieve near-perfect accuracy (99.2--100\%) at classifying epistemic versus non-epistemic statements from middle layers onward, with high selectivity scores (0.45--0.49) confirming genuine encoding; (2) non-epistemic agreement is more linearly decodable (87--95\%) than factual truth (54--59\% across mixed domains), suggesting LLMs encode sentiment and normativity more strongly than truth value; and (3) belief encoding follows consistent layer-wise patterns across model scales, with epistemic truth peaking in middle layers and type classification emerging robustly by normalized depth 0.50--0.67.
These findings reveal that LLMs develop distinct representational subspaces for different belief categories, with implications for hallucination detection and targeted belief interventions.
